\section{Supplementary Theory and Full Proofs}

\subsection{Operational Gap Notation}

For reference, the experimental gap can be written without adding another
formal theorem hypothesis. Let:
\begin{itemize}[leftmargin=1.5em]
    \item $L_B(C)$ be exact lowering into target basis $B$,
    \item $\kappa$ be a conservative structural cost,
    \item $\mathcal{T}_K(C)$ be the bounded set of semantic candidates produced
    by the semantic lift under budget $K$,
    \item $\mathfrak{A}$ be a configured post-lowering pipeline family used for
    an experiment.
\end{itemize}

Define the $(K,w)$-representation gap of $C$ by
\[
\mathrm{Gap}^{K,w}_{\mathrm{rep}}(C)
:=
\inf_{A \in \mathfrak{A}} \kappa(A(L_B(C)))
-
\inf_{T \in \mathcal{T}_K(C)} \kappa(L_B(T)).
\]
Define the operational recoverability gap by
\[
\mathrm{Gap}^{K,w}_{\mathrm{rec}}(C)
:=
\kappa(C_{\mathrm{flat}}^\star) - \kappa(C_{\mathrm{term}}^\star),
\]
where $C_{\mathrm{flat}}^\star$ is the best candidate recovered from the flat
basis-lowered representation and $C_{\mathrm{term}}^\star$ is the best
candidate recovered from the bounded semantic term representation.

The meaning is simple: semantics may survive lowering, while cheap
recoverability of the opportunity may not. This notation is operational and
experimental; the formal lower bound in the main text is the budget-explicit
final-output theorem for finite-window transducers.

\begin{definition}[Aggregate recovery]
\label{def:aggregation-equivalent-recovery}
Consider the Fourier-layer family
$C_{m,r}=H_{\Lambda}D_m(\boldsymbol{\theta})^rH_{\Lambda}$ with
$D_m(\boldsymbol{\theta})=\prod_{j=1}^m e^{-i\theta_jP_j/2}$. A compiler
performs \emph{aggregation-equivalent recovery} on this family if, for each
non-removable term $P_j$, its successful output can be associated with an
$O(1)$-size record, subcircuit, or algebraic object whose coefficient is
equivalent modulo $2\pi$ to the aggregate $r\theta_j$. The object may be
explicit, as in a phase-polynomial or commuting-diagonal IR, or implicit, as
in a specialized rewrite, lookup, ZX transformation, or synthesis routine that
computes the same aggregate diagonal action without storing all $r$ lowered
occurrences separately.
\end{definition}

\begin{proposition}[Operational recovery channels]
\label{prop:classification-boundary}
For the Fourier-layer family, a pipeline that recovers the constant
fixed-width output can be classified by the channel through which it carries
the aggregate coefficient information. Examples include:
\begin{enumerate}[leftmargin=1.5em]
    \item store enough persistent summary information to carry the aggregate;
    \item lift enough lowered volume into a richer IR to reconstruct the
    aggregate;
    \item use nonlocal provenance information from before basis lowering that
    already contains the aggregate;
    \item reconstruct a global semantic representation equivalent to a
    phase-polynomial, ZX, stabilizer-phase, or commuting-diagonal Hamiltonian
    object.
\end{enumerate}
Consequently, a successful future flat-looking implementation would not by
itself refute the main theorem. It would instead classify the implementation by
identifying where the aggregate information crosses the finite-state boundary,
or by showing that the formal transducer model is missing a cheaper channel.
\end{proposition}

\begin{proof}[Justification]
The fixed-width theorem lower-bounds final output length for constant-budget
finite-window transducers on a certified noncolliding family. A pipeline that
achieves the semantic constant output must still represent the aggregate
coefficients that distinguish the noncolliding cases. The listed channels are
implementation-level ways to carry that information. This is a classification
statement rather than an optimality claim.
\end{proof}

\subsection*{Additional Structural Families}

The artifact also targets phase-ladder, mirrored, and conjugation patterns that
appear in QPE/QFT-style workloads and Grover-style workloads. These families
are not used here as additional theorem instances. They are best read as
structural stress families: the semantic controller preserves repeated
phase-heavy ladders, mirrored shells, and conjugation shells long enough to
compare projected costs before materializing the selected candidate.

For phase-ladder workloads, the relevant empirical question is whether bounded
semantic terms preserve a repeated commuting-diagonal opportunity that is
harder to recover after early lowering. For mirrored and conjugation workloads,
the relevant empirical question is whether the artifact recognizes repeated
shell structure and avoids materializing every candidate before selection. The
Grover and mirrored-shell experiments in Appendix~C are therefore empirical
positive cases and robustness checks, not support for a separate theorem-level
family separation.

\subsection{Self-Contained Rational CP Period Certificate}
\label{app:sec:rational-period}

The canonical four-qubit rational witness has diagonal block
\begin{equation*}
\begin{aligned}
D={}&\prod_{j=0}^{3}RZ_j(\theta_j)
 \prod_{0\leq a<b\leq3}CP_{ab}(\theta_{ab}),\\
\boldsymbol\theta_{RZ}/\pi
={}&(1/7,1/8,1/9,1/10),\\
(\theta_{01},\theta_{02},\theta_{03})/\pi
={}&(1/12,1/13,1/14),\\
(\theta_{12},\theta_{13},\theta_{23})/\pi
={}&(1/14,1/15,1/16).
\end{aligned}
\end{equation*}
For a reduced rational angle $\theta/\pi=p/q$, both $RZ(\theta)$ and
$CP(\theta)$ have projective single-term order
\begin{equation*}
 \operatorname{ord}(p/q)
 =\frac{2q}{\gcd(|p|,2q)},
\end{equation*}
because the relevant relative phase is $p/q$ modulo $2$.  The complete
single-term certificate is
\begin{center}
\begin{tabular}{lccc}
\toprule
term & support & $\theta/\pi$ & $T_j$\\
\midrule
$RZ_0$ & $0$ & $1/7$ & $14$\\
$RZ_1$ & $1$ & $1/8$ & $16$\\
$RZ_2$ & $2$ & $1/9$ & $18$\\
$RZ_3$ & $3$ & $1/10$ & $20$\\
$CP_{01}$ & $01$ & $1/12$ & $24$\\
$CP_{02}$ & $02$ & $1/13$ & $26$\\
$CP_{03}$ & $03$ & $1/14$ & $28$\\
$CP_{12}$ & $12$ & $1/14$ & $28$\\
$CP_{13}$ & $13$ & $1/15$ & $30$\\
$CP_{23}$ & $23$ & $1/16$ & $32$\\
\bottomrule
\end{tabular}
\end{center}
Hence
\begin{equation*}
\begin{aligned}
T_{\mathrm{terms}}
&=\operatorname{lcm}(14,16,18,20,24,26,28,28,30,32)\\
&=2^5\cdot3^2\cdot5\cdot7\cdot13=131040,
\end{aligned}
\end{equation*}
which proves that $D^{131040}$ is a global phase.  To certify minimality
for the product rather than merely obtain this common upper bound, subtract
the phase of $0000$.  For $z=z_0z_1z_2z_3$, the exact relative phase is
\begin{equation*}
\begin{aligned}
w(z)={}&z_0/7+z_1/8+z_2/9+z_3/10\\
 &+z_0z_1/12+z_0z_2/13+z_0z_3/14\\
 &+z_1z_2/14+z_1z_3/15+z_2z_3/16
 \pmod 2.
\end{aligned}
\end{equation*}
In particular,
\begin{equation*}
\begin{aligned}
 w(1011)&=\frac{37007}{65520},\\
 \operatorname{ord}(w(1011))
 &=\frac{2\cdot65520}{\gcd(37007,2\cdot65520)}\\
 &=131040.
\end{aligned}
\end{equation*}
Thus no smaller positive power of $D$ is a global phase, and
$T=131040$ exactly.  It follows that
$D^1,\ldots,D^R$ are pairwise distinct up to global phase for every
$R<T$.

The machine certificate uses only Python \texttt{fractions.Fraction} and
prints all 16 relative numerators, denominators, and orders.  The frozen
integer output contains
\begin{equation*}
\begin{gathered}
\texttt{term\_order\_lcm=131040},\\
\texttt{relative\_order\_lcm=131040},\\
\texttt{minimality\_state\_1011\_order=131040},\\
\texttt{exact\_global\_phase\_period\_T=131040}.
\end{gathered}
\end{equation*}
The reproducibility files and hashes are listed below; each digest is split
at its 32-hex midpoint only for typesetting.
\begin{itemize}[leftmargin=1.5em]
 \item verifier: \path{research/rational_witness/verify_period.py},
 SHA256 \texttt{c7c90f26e00e515383e94e4bbd707574}\allowbreak
 \texttt{eb3aabf94ad71c3cd5e2cdb8d5db7e76};
 \item integer output:
 \path{research/rational_witness/period_certificate.txt}, SHA256
 \texttt{375295d776f21ea2a82ee4f2bab5205d}\allowbreak
 \texttt{25bfef85ee864900baacfbd80be9bd6e};
 \item manifest: \path{research/rational_witness/SHA256SUMS}.
\end{itemize}
From the repaired-manuscript root, the byte-for-byte check is
\begin{quote}
\footnotesize
\texttt{cd research/rational\_witness}\\
\texttt{python3 verify\_period.py \textbackslash}\\
\texttt{\quad --check period\_certificate.txt}
\end{quote}

\subsection{Full Proofs for the Transducer Lower Bound}
\label{app:sec:full-proofs}

This section supplies the full finite-state proof layer underlying
\Cref{thm:final-output,cor:finite-tail,cor:counting-toll}. The main text states
the shorter version; the appendix keeps the pumpability data explicit.
The following statements restate the main-text theorem and corollaries with
full hypotheses and complete proofs; they are not additional independent
results.

\subsubsection{Pumpable Pair Families}

Persistent store can depend on $r$, so the induction tracks both a store value
and a tape word. In this subsection, $D$ in pumpability data denotes a period
modulus, not the diagonal witness layer in $C_r=H_\Lambda D^rH_\Lambda^{-1}$.
The next definition restates the main-text pumpable-pair definition with
explicit residue data for the full proof.

\begin{definition}[Pumpable pair family, restated with explicit residue data]
\label{app:def:pumpable}
Let $\Gamma$ be a finite alphabet and $P$ a finite set of store values. A family
\[
  X(r)=(\beta(r),W(r))\in P\times\Gamma^*,\qquad r\in\mathbb{N},
\]
is \emph{pumpable} if there exist integers $h\ge0$ and $D\ge1$ such that for
every residue $\rho\in\{0,\ldots,D-1\}$ there are
\[
  \beta_\rho\in P,
  \qquad
  A_\rho,B_\rho,C_\rho\in\Gamma^*,
\]
with
\[
  X(h+\rho+kD)=(\beta_\rho,A_\rho B_\rho^k C_\rho)
\]
for every $k\ge0$. The word $B_\rho$ is the pump word on residue $\rho$. The
residue is \emph{live} if $B_\rho\ne\epsilon$ and \emph{collapsed} if
$B_\rho=\epsilon$.
\end{definition}

\begin{lemma}[Input pumpability]
\label{app:lem:input-pumpable}
For the fixed-width witness $T_0(r)=\Pi M^r\Sigma$, the pair family consisting
of the initial persistent store and $T_0(r)$ is pumpable with $h=0$, $D=1$,
pump word $M$, and fixed store value.
\end{lemma}

\begin{proof}
Take $A_0=\Pi$, $B_0=M$, and $C_0=\Sigma$.
\end{proof}

\subsubsection{One-Pass Propagation}

Fix a precision bound $q$ and a charged alphabet $\Gamma=\Gamma_B(q)$. Under
this bound, a pass of $\mathcal{A}$ can be viewed as an ordinary deterministic
subsequential transducer over a finite alphabet. Its finite state consists of
the ephemeral state, the persistent store, and the $w$-token window. The
end-of-pass persistent store is one component of the final finite state.

\begin{lemma}[One-pass pumpability propagation]
\label{app:lem:one-pass-pump}
Let $X(r)=(\beta(r),W(r))$ be a pumpable pair family over a finite input
alphabet. Let $\tau$ be one deterministic finite-window, finite-burst pass whose
output alphabet and store set are finite. Let
\[
  \tau(\beta(r),W(r))=(\beta'(r),W'(r))
\]
be the resulting final store and output tape. Then
$X'(r)=(\beta'(r),W'(r))$ is pumpable.
\end{lemma}

\begin{proof}
Fix pumpability data
$h,D,\{\beta_\rho,A_\rho,B_\rho,C_\rho\}_\rho$ for $X$. We first treat one
residue $\rho$, then recombine all residues into $r$-indexed pumpability data
for $X'$.

For a fixed residue $\rho$, write
\[
  W_k=A B^k C,
  \qquad
  \beta_k=\beta,
\]
where $k\ge0$ and $A,B,C,\beta$ abbreviate the residue data. Absorb the
ephemeral state, persistent store, and window into the pass's finite state set
$S$. Starting from the state determined by $\beta$, process $A$. This emits a
fixed word $O_A$ and reaches a state $s_A\in S$.

Let $f:S\to S$ be the state map obtained by reading the block $B$, and let
$o(s)\in(\Gamma')^*$ be the output emitted while reading $B$ starting in state
$s$. The sequence
\[
  s_A,\ f(s_A),\ f^2(s_A),\ldots
\]
in the finite set $S$ is ultimately periodic. Thus there exist $u\ge0$ and
$\ell\ge1$ with $u+\ell\le |S|+1$ such that
\[
  f^{t+\ell}(s_A)=f^t(s_A)
  \qquad\text{for all }t\ge u.
\]
For $k=u+v\ell+\eta$ with $0\le\eta<\ell$, the output while reading $B^k$ is
\[
\begin{aligned}
  &\bigl(o(s_A)o(f(s_A))\cdots o(f^{u-1}(s_A))\bigr)\\
  &\quad\cdot\Bigl(o(f^u(s_A))\cdots
  o(f^{u+\ell-1}(s_A))\Bigr)^v\\
  &\quad\cdot\bigl(o(f^u(s_A))\cdots
  o(f^{u+\eta-1}(s_A))\bigr).
\end{aligned}
\]
After this, the pass reads the fixed suffix $C$, emits a word depending only
on the state $f^{u+\eta}(s_A)$, and ends in a persistent store value depending
only on the same state.

Restore the residue subscripts. For each original residue
$\rho\in\{0,\ldots,D-1\}$, let $u_\rho\ge0$ and $\ell_\rho\ge1$ be the
preperiod and cycle length obtained above, let $f_\rho$, $s_{A_\rho}$, and
$O_{A_\rho}$ be the corresponding state map, post-prefix state, and prefix
output. For $0\le\eta<\ell_\rho$, let $F_{\rho,\eta}$ be the complete word
emitted while reading the fixed suffix $C_\rho$ from state
$f_\rho^{u_\rho+\eta}(s_{A_\rho})$, including the end-of-tape finalizer output,
and define
\begin{align*}
  P_\rho &:=
  o\bigl(s_{A_\rho}\bigr)o\bigl(f_\rho(s_{A_\rho})\bigr)\cdots
  o\bigl(f_\rho^{u_\rho-1}(s_{A_\rho})\bigr),\\
  Q_\rho &:=
  o\bigl(f_\rho^{u_\rho}(s_{A_\rho})\bigr)\cdots
  o\bigl(f_\rho^{u_\rho+\ell_\rho-1}(s_{A_\rho})\bigr),\\
  R_{\rho,\eta} &:=
  o\bigl(f_\rho^{u_\rho}(s_{A_\rho})\bigr)\cdots
  o\bigl(f_\rho^{u_\rho+\eta-1}(s_{A_\rho})\bigr)\cdot F_{\rho,\eta}.
\end{align*}
Let $\beta'_{\rho,\eta}$ be the end-of-pass persistent store value reached from
that state. Define
\[
\begin{aligned}
  u^*&:=\max_{\rho} u_\rho,
  &\ell^*&:=\lcmop_{\rho}\ell_\rho,\\
  h'&:=h+D u^*,
  &D'&:=D\ell^*.
\end{aligned}
\]
Each new residue $\rho'\in\{0,\ldots,D'-1\}$ decomposes uniquely as
\[
  \rho'=\rho+D\tilde\eta,
  \qquad
  \rho\in\{0,\ldots,D-1\},\quad
  \tilde\eta\in\{0,\ldots,\ell^*-1\}.
\]
For $r=h'+\rho'+k'D'$ with $k'\ge0$, the old representation
$r=h+\rho+kD$ has
\[
  k=u^*+\tilde\eta+k'\ell^*.
\]
Since $u^*\ge u_\rho$, we have $k\ge u_\rho$ for every $k'\ge0$; and since
$\ell_\rho\mid\ell^*$, the residue of $k-u_\rho$ modulo $\ell_\rho$ does not
depend on $k'$. Set
\[
  \eta:=(u^*+\tilde\eta-u_\rho)\bmod \ell_\rho,
  \qquad
  v_0:=\Bigl\lfloor\frac{u^*+\tilde\eta-u_\rho}{\ell_\rho}\Bigr\rfloor\ge0,
\]
so that
\[
  k=u_\rho+(v_0+k'\ell^*/\ell_\rho)\ell_\rho+\eta.
\]
By the output decomposition above, for every $k'\ge0$,
\[
  X'(h'+\rho'+k'D')
  =\bigl(\beta'_{\rho,\eta},\ A'_{\rho'}(B'_{\rho'})^{k'}C'_{\rho'}\bigr),
\]
where
\[
  A'_{\rho'}:=O_{A_\rho}P_\rho Q_\rho^{v_0},
  \qquad
  B'_{\rho'}:=Q_\rho^{\ell^*/\ell_\rho},
  \qquad
  C'_{\rho'}:=R_{\rho,\eta}.
\]
This is exactly pumpability of $X'$ with data $(h',D')$ and residue data
$\{\beta'_{\rho,\eta},A'_{\rho'},B'_{\rho'},C'_{\rho'}\}_{\rho'}$.
\end{proof}

\begin{corollary}[Per-tape induction, for proof]
\label{app:cor:per-tape}
Assume $\mathcal{A}$ has constant budget, so all executions use one finite
alphabet $\Gamma_B(q_*)$. For every pass index $i=0,1,\ldots,p$, the pair
family consisting of the persistent store after pass $i$ and the tape $T_i(r)$
is pumpable.
\end{corollary}

\begin{proof}
The base case is \Cref{app:lem:input-pumpable}. Apply
\Cref{app:lem:one-pass-pump} successively to passes $1,\ldots,p$.
\end{proof}

\begin{remark}[Unary scratch tapes]
\label{app:rem:unary}
If pass 1 writes one marker per input period, then
$T_1(r)=\mathtt{marker}^r$ is pumpable with live pump word
$\mathtt{marker}$. \Cref{app:lem:one-pass-pump} applies to every later
pass. A later constant-budget pass may output one marker every $c$ input
markers, permute markers, or add a fixed periodic decoration, but it remains
pumpable. It cannot turn $\mathtt{marker}^r$ into a family of $O(1)$ many
correct aggregate circuits for all noncolliding $r$; if the final pump
collapses, the final output repeats on a residue class.
\end{remark}

\subsubsection{Collapse Versus Final-Output Length}

\begin{lemma}[Pumpable outputs are bounded-or-linear on each residue]
\label{app:lem:bounded-or-linear}
Let $W(r)$ be a pumpable word family with data
$(h,D,A_\rho,B_\rho,C_\rho)$. For $r=h+\rho+kD$:
\begin{enumerate}[label=(\alph*)]
\item if $B_\rho=\epsilon$, then $W(h+\rho+kD)$ is independent of $k$;
\item if $B_\rho\ne\epsilon$, then
\[
  |W(h+\rho+kD)|\ge k.
\]
\end{enumerate}
Consequently, on a live residue,
\[
  |W(r)|\ge \frac{r-h-D+1}{D}
\]
for all $r\ge h+D$ in that residue.
\end{lemma}

\begin{proof}
Both statements are immediate from
$W(h+\rho+kD)=A_\rho B_\rho^k C_\rho$. If $B_\rho\ne\epsilon$, then
$|B_\rho|\ge1$, so $|A_\rho B_\rho^k C_\rho|\ge k$. Since
$r=h+\rho+kD$ and $0\le\rho<D$, $k\ge(r-h-D+1)/D$.
\end{proof}

\begin{lemma}[Correct noncolliding final outputs cannot collapse]
\label{app:lem:no-collapse}
Assume finite-range noncollision (Assumption~\ref{ass:noncollision}) with
$N(R)=R$ on the interval $1\le r\le R$. Suppose the final output family
$T_p(r)$ is pumpable with data $(h,D,A_\rho,B_\rho,C_\rho)$ and
$\mathcal{A}$ is correct on every $r\in\{1,\ldots,R\}$. Without loss of
generality take $h\ge1$: if $h=0$, replace $h$ by $h+D$; this preserves
pumpability after discarding the finite prefix of indices $k=0$ on each residue
and absorbing one copy of $B_\rho$ into $A_\rho$. Then every residue class
appearing twice in the tail $\{h,h+1,\ldots,R\}$ is live. In particular, for
every $r\in\{h+D,\ldots,R\}$,
\[
  |T_p(r)|\ge \frac{r-h-D+1}{D}.
\]
\end{lemma}

\begin{proof}
Let $r=h+\rho+kD$ with $k\ge1$ and $r\le R$. Then
$r-D=h+\rho+(k-1)D$ also lies in the tail, and by the convention $h\ge1$ it
satisfies $r-D\ge h\ge1$, so both $r$ and $r-D$ lie in the correctness domain
$\{1,\ldots,R\}$. If $B_\rho=\epsilon$,
\Cref{app:lem:bounded-or-linear} gives
\[
  T_p(r)=T_p(r-D)
\]
as literal circuit strings. Hence the two outputs have the same unitary.
Correctness at both $r$ and $r-D$ would imply
\[
U(\mathcal{L}_B(C_r))=U(\mathcal{L}_B(C_{r-D}))
\]
up to global phase, which by exactness of the lowering means $D^r\equiv
D^{r-D}$ up to global phase, contradicting Assumption~\ref{ass:noncollision}
with $N(R)=R$. Therefore $B_\rho\ne\epsilon$, and
\Cref{app:lem:bounded-or-linear} gives the stated length bound.
\end{proof}

\subsubsection{Final-Output Theorem and Corollaries}

\begin{theorem}[$p$-pass final-output lower bound, restated with full hypotheses]
\label{app:thm:final-output}
Fix a fixed-width diagonal witness family
$C_r=H_\Lambda D^rH_\Lambda^{-1}$ with exact lowering
$\mathcal{L}_B(C_r)=\Pi M^r\Sigma$, and assume unbounded noncollision
(Assumption~\ref{ass:unbounded-noncollision}): the powers $D^r$ are pairwise
distinct up to global phase for all $r\ge1$, so $N(R)=R$ on every range. Let
$\mathcal{A}$ be a uniform deterministic $p$-pass finite-window, finite-burst
transducer with constant budget that is correct on the input
$\mathcal{L}_B(C_r)=\Pi M^r\Sigma$ for every $r\ge1$.
Here uniform means one fixed deterministic transducer, with the same passes,
transition rules, finite state, window, burst bound, and store bound for all
repetition counts $r$.

Then there exist constants $h\ge1$ and $D\ge1$, depending on $\mathcal{A}$ and
the fixed witness but not on $r$, such that for every $r\ge h+D$,
\[
  |T_p(r)|\ge \frac{r-h-D+1}{D}=\Omega_{\mathcal{A}}(r).
\]
Equivalently, no constant-budget transducer in this model that is correct on
the whole noncolliding fixed-width family can emit $o(r)$-length final outputs
on that family.
\end{theorem}

\begin{proof}
By constant budget, all executions use a finite charged token alphabet.
\Cref{app:cor:per-tape} gives pumpability of the final pair family and
hence of $T_p(r)$, with data $(h,D)$ depending only on $\mathcal{A}$ and the
witness. By the convention of \Cref{app:lem:no-collapse}, take $h\ge1$,
so every comparison point $r-D\ge h\ge1$ used below lies in the correctness
domain.

Fix any $r\ge h+D$ and any $R\ge r$. By Assumption~\ref{ass:unbounded-noncollision},
$N(R)=R$, and $\mathcal{A}$ is correct on every $r'\in\{1,\ldots,R\}$. Thus
\Cref{app:lem:no-collapse} applies and gives the displayed bound at $r$.
In particular every final residue is live: if any final residue collapsed, two
sufficiently large inputs in that residue would receive the same final output
string, contradicting correctness and noncollision.
\end{proof}

\paragraph{Correctness quantifier.}
Correctness in \Cref{app:thm:final-output} is required at every $r\ge1$; the
conclusion then holds at every $r\ge h+D$, without exception. This global
quantifier is not decorative: the collapse argument in
\Cref{app:lem:no-collapse} compares the outputs at $r$ and $r-D$ and
needs correctness at both points. Correctness at isolated points is not
sufficient; if $\mathcal{A}$ were correct at $r$ but not at $r-D$, identical
final strings at those two inputs would yield no contradiction.

\begin{corollary}[Finite-range tail form, restated with full hypotheses]
\label{app:cor:finite-tail}
Let $\mathcal{A}$ be a constant-budget transducer, and let $(h,D)$ be
pumpability data for its final pair family as provided by
\Cref{app:cor:per-tape}. Suppose finite-range noncollision
(Assumption~\ref{ass:noncollision}) holds with $N(R)=R$ on the range
$1\le r\le R$ and $\mathcal{A}$ is correct on every
$r\in\{1,\ldots,R\}$. Then
\[
  |T_p(r)|\ge \frac{r-h-D+1}{D}
  \qquad\text{for every }h+D\le r\le R.
\]
The constants $(h,D)$ are determined by $\mathcal{A}$ and the witness family
alone; in particular they do not depend on $R$. No statement is made for
transducers outside the constant-budget regime: there, pumpability data in the
sense of \Cref{app:def:pumpable} need not exist, and any finite-range
inequality would require explicit control of pumping constants not supplied by
this theorem.
\end{corollary}

\begin{proof}
Constant budget puts all executions over one finite charged alphabet, so
\Cref{app:cor:per-tape} supplies pumpability data $(h,D)$ for the
final tape. The construction of this data uses only the fixed transducer, its
finite alphabets and state sets, and the fixed witness words $\Pi,M,\Sigma$;
it is obtained before any finite endpoint $R$ is chosen. Hence $(h,D)$ are
independent of $R$. \Cref{app:lem:no-collapse} then applies verbatim on
the range $1\le r\le R$.
\end{proof}

\begin{corollary}[Bit-level output counting toll, restated for proof]
\label{app:cor:counting-toll}
Assume finite-range noncollision supplies $N(R)$ pairwise distinct target
unitaries, up to global phase, on a range. Suppose a deterministic exact
compiler emits final outputs $T_p(r)$. Then
\[
  \max_{1\le r\le R}B_{\mathrm{self}}(T_p(r))
  \ge \left\lceil\log_2 N(R)\right\rceil.
\]
The same inequality holds for a fixed family-relative code when one public
dictionary is shared by the whole range.  For the rational-angle CP witness,
only the displayed per-range inequality is claimed on fixed ranges
$R<T=131040$.
\end{corollary}

\begin{proof}
A circuit string determines its unitary, so a deterministic exact compiler
that is correct on two inputs whose target unitaries are distinct up to global
phase must emit two distinct output strings. The noncollision hypothesis
therefore forces at least $N(R)$ pairwise distinct codewords.  If their
prefix-free lengths are $\ell_1,\ldots,\ell_{N(R)}$, Kraft's inequality
gives $\sum_i2^{-\ell_i}\leq1$, so
$N(R)2^{-\max_i\ell_i}\leq1$.  Hence the largest length is at least
$\lceil\log_2N(R)\rceil$.  The argument applies independently to either
fixed code and does not convert token count or gate count into bits.
\end{proof}

\subsection{Bookkeeping for the Masked-Share Cut}
\label{app:sec:streaming-proof}

This subsection records the restart property used by the masked-share
tradeoff (\cref{thm:main-tradeoff}).  It also makes explicit why a
materialized update log, phase table, graph, RAM/cache layout, old-input
spool, dynamic-pass state, or random seed is part of the cut message rather
than an uncharged reconstruction channel.

\begin{lemma}[Restart from the serialized cut]
\label{app:lem:reconstruction}
Fix $\mathcal A$ as in \cref{def:compiler}.  At the inter-round cut
$c_{\mathrm{MS}}$, the five fields in \eqref{eq:info-budget}, together
with any prefix already committed to the final write-once output, determine
the continuation of that run for every supplied second-round continuation
$v$.  For a randomized run this statement is conditional on the serialized
seed/tape and cursor.
\end{lemma}

\begin{proof}
The control field restores the exact transition, all head positions,
dynamic pass/termination schedule, and random seed/tape cursor.  The store
and window fields restore RAM, stack, hash/cache metadata, nonparameter
machine data, and ordered live token structures.  The parameter field
restores every omitted coefficient by its location tag.  The IR field
restores every other live intermediate tape, DAG, table, addressed record,
and any consumed input block retained for rescanning.  These disjoint fields
therefore reconstruct the complete machine configuration.  Conditional on
the serialized random tape, the receiver can append the supplied second
round and execute the same dynamically chosen continuation.  The separately
transmitted final-output prefix is prepended to the suffix produced after
restart.  No additional state, old-input file, or public-family side channel
is used.
\end{proof}
